\documentclass[11pt,a4paper]{article}
\usepackage[margin=19mm,headheight=15pt]{geometry}
\usepackage{fontspec}
\setmainfont{Noto Serif}
\usepackage{amsmath,amssymb,mathtools}
\usepackage{xcolor}
\usepackage{fancyhdr}
\usepackage{enumitem}
\usepackage{microtype}
\setlength{\parindent}{0pt}
\setlength{\parskip}{6pt}
\definecolor{navy}{HTML}{183153}
\definecolor{lightgray}{HTML}{F3F5F7}
\pagestyle{fancy}\fancyhf{}
\lhead{\small ME 5309 -- Mid-Semester Examination}
\rhead{\small ME 5309}
\cfoot{\thepage}
\begin{document}
\begin{center}
{\Large\bfseries ME 5309 -- Mid-Semester Examination}\\[3pt]
{\large Answer Booklet}\\[8pt]
\end{center}
\textbf{Instructions:} Attempt any three questions. Each question carries 10 marks. Show working where appropriate.
\section*{Question 2 \hfill [10 marks]}
\subsection*{Q2(a) \hfill [2 marks]}

ẋ₁ = x₂ and ẋ₂ = (τ-b x₂+mgℓ sin θ)/J.
At equilibrium:\[\begin{aligned}τ_0 = -mg\ell sin \theta_e,\\ \text{an equilibrium exists for |τ₀| less than or equal to 2.4525 N m}\end{aligned}\]

\subsection*{Q2(b) \hfill [2 marks]}

The Jacobian has A₁₂ = 1 and A₂₂ = -b/J = -2.5.
The input matrix is B = [0, 50]ᵀ.

\subsection*{Q2(c) \hfill [4 marks]}

At steady state the derivative terms vanish.
The proportional torque is Kp(θd-θ), which must balance mgℓ sin θ.

\subsection*{Q2(d) \hfill [2 marks]}

Gravity compensation makes the steady-state error zero.
It does not change the closed-loop poles because it only adds a constant torque.

\section*{Question 3 \hfill [10 marks]}
\subsection*{Q3(a) \hfill [2 marks]}

\vspace{2mm}

\subsection*{Q3(b) \hfill [2 marks]}

\vspace{2mm}

\subsection*{Q3(c) \hfill [4 marks]}

A-LC=[[-ℓ₁,1], [-ℓ₂,-2]].
I obtain p(s)=s²+(ℓ₁+2)s+(ℓ₁+ℓ₂).
Matching (s+10)² gives ℓ₁=18 and ℓ₂=82.

\subsection*{Q3(d) \hfill [2 marks]}

An observer can estimate velocity from position measurements.
A filter can be safer because it smooths noise.

\section*{Question 4 \hfill [10 marks]}
\subsection*{Q4(a) \hfill [2 marks]}

The characteristic polynomial is s²+4s+8, so I report poles -4 and -8.
Unity DC gain requires Nff=8.

\subsection*{Q4(b) \hfill [2 marks]}

The nominal controller has unity DC gain, so θss=1 and the tracking error is zero even after the actuator gain changes.

\subsection*{Q4(c) \hfill [4 marks]}

\[\begin{aligned}k_1=18, k_2=6, K_I=32\end{aligned}\]

\subsection*{Q4(d) \hfill [2 marks]}

Integral action gives zero steady-state error if stable.
A practical problem is integrator windup.

\end{document}